
\documentclass{article}
\usepackage{colortbl}
\usepackage{makecell}
\usepackage{multirow}
\usepackage{supertabular}

\begin{document}

\newcounter{utterance}

\centering \large Interaction Transcript for game `cladder', experiment `full\_v1.5\_default', episode 8837 with selfrefine{-}llama{-}3.1{-}8b{-}instruct.
\vspace{24pt}

{ \footnotesize  \setcounter{utterance}{1}
\setlength{\tabcolsep}{0pt}
\begin{supertabular}{c@{$\;$}|p{.15\linewidth}@{}p{.15\linewidth}p{.15\linewidth}p{.15\linewidth}p{.15\linewidth}p{.15\linewidth}}
    \# & \multicolumn{2}{c}{Player} && \multicolumn{2}{c}{Game Master} \\
    \hline

    \theutterance \stepcounter{utterance}  
    & & & \multicolumn{4}{p{0.6\linewidth}}{
        \cellcolor[rgb]{0.9,0.9,0.9}{
            \makecell[{{p{\linewidth}}}]{
                \texttt{\tiny{[P1$\langle$GM]}}
                \texttt{You are an expert in causal inference. The following question is not a typical commonsense query, but rather a meticulously designed question created by a professor specializing in causal inference, intended to assess the students' mastery of the course content.} \\
\texttt{Imagine a self{-}contained, hypothetical world with only the following conditions, and without any unmentioned factors or causal relationships: Yupt has a direct effect on zupj and muvq. Xyfo has a direct effect on zupj and muvq. Zupj has a direct effect on muvq. Xyfo is unobserved. For those who are not yupt and are not zupj, the probability of muvq is 44\%. For those who are not yupt and are zupj, the probability of muvq is 51\%. For those who are yupt and are not zupj, the probability of muvq is 57\%. For those who are yupt and are zupj, the probability of muvq is 10\%. For those who are not yupt and are not xyfo, the probability of zupj is 42\%. For those who are not yupt and are xyfo, the probability of zupj is 87\%. For those who are yupt and are not xyfo, the probability of zupj is 26\%. For those who are yupt and are xyfo, the probability of zupj is 44\%. The overall probability of xyfo is 22\%. Does yupt negatively affect muvq through zupj?} \\
\texttt{Start your answer with 'Yes' or 'No', followed by additional reasoning or evidence to support your explanation.} \\
            }
        }
    }
     \\ \\

    \theutterance \stepcounter{utterance}  
    & & \multicolumn{4}{p{0.6\linewidth}}{
        \cellcolor[rgb]{0.95,0.95,0.95}{
            \makecell[{{p{\linewidth}}}]{
                \texttt{\tiny{[P1$\leftrightarrow$P]}}
                \texttt{{[}DRAFT{]}} \\
\texttt{No, yupt does not negatively affect muvq through zupj.} \\
\\ 
\texttt{To assess the causal effect of yupt on muvq through zupj, we can use the backdoor criterion, which involves identifying a path from yupt to muvq that is unconfounded by any other variables. } \\
\\ 
\texttt{First, we need to identify the path from yupt to muvq through zupj. The path is: yupt → zupj → muvq.} \\
\\ 
\texttt{Next, we need to check if this path is unconfounded by any other variables. We can do this by checking if the probability of muvq is the same for those who are yupt and zupj as it is for those who are not yupt and zupj.} \\
\\ 
\texttt{However, we are given that for those who are yupt and are zupj, the probability of muvq is 10\%, and for those who are not yupt and are zupj, the probability of muvq is 51\%. Since these probabilities are not the same, the path from yupt to muvq through zupj is confounded by the variable xyfo.} \\
\\ 
\texttt{However, we are given that the overall probability of xyfo is 22\%, and we can use this information to calculate the probability of muvq for those who are yupt and zupj, given that they are not xyfo. We can do this by using the law of total probability:} \\
\\ 
\texttt{P(muvq|yupt, zupj) = P(muvq|yupt, zupj, xyfo) * P(xyfo|yupt, zupj) + P(muvq|yupt, zupj, not xyfo) * P(not xyfo|yupt, zupj)} \\
\\ 
\texttt{We are given that P(muvq|yupt, zupj, xyfo) = 10\% and P(xyfo|yupt, zupj) = 44\% (since the probability of xyfo is 22\% and the probability of zupj is 44\% for those who are yupt). We are also given that P(muvq|yupt, zupj, not xyfo) = 57\% (since the probability of muvq is 57\% for those who are yupt and not zupj, and the probability of zupj is 26\% for those who are yupt and not xyfo). We can also calculate P(not xyfo|yupt, zupj) = 1 {-} P(xyfo|yupt, zupj) = 1 {-} 44\% = 56\%.} \\
\\ 
\texttt{Plugging in these values, we get:} \\
\\ 
\texttt{P(muvq|yupt, zupj) = 10\% * 44\% + 57\% * 56\% ≈ 32.28\%} \\
\\ 
\texttt{However, we are given} \\
            }
        }
    }
    & \\ \\

    \theutterance \stepcounter{utterance}  
    & & \multicolumn{4}{p{0.6\linewidth}}{
        \cellcolor[rgb]{0.95,0.95,0.95}{
            \makecell[{{p{\linewidth}}}]{
                \texttt{\tiny{[P1$\leftrightarrow$P]}}
                \texttt{{[}SELF{-}FEEDBACK{]}} \\
\texttt{FORMAT: ok} \\
\texttt{RULES: ok} \\
\texttt{STRATEGY: problem} \\
\texttt{REASON: The answer incorrectly concludes that yupt does not negatively affect muvq through zupj, despite the fact that the probability of muvq is lower for those who are yupt and zupj compared to those who are not yupt and zupj, indicating a negative effect.} \\
            }
        }
    }
    & \\ \\

    \theutterance \stepcounter{utterance}  
    & \multicolumn{4}{p{0.6\linewidth}}{
        \cellcolor[rgb]{0.561,0.737,0.561}{
            \makecell[{{p{\linewidth}}}]{
                \texttt{\tiny{[P1$\rangle$GM]}}
                \texttt{No, yupt does not negatively affect muvq through zupj.} \\
\\ 
\texttt{To assess the causal effect of yupt on muvq through zupj, we can use the backdoor criterion, which involves identifying a path from yupt to muvq that is unconfounded by any other variables.} \\
\\ 
\texttt{First, we need to identify the path from yupt to muvq through zupj. The path is: yupt → zupj → muvq.} \\
\\ 
\texttt{Next, we need to check if this path is unconfounded by any other variables. We can do this by checking if the probability of muvq is the same for those who are yupt and zupj as it is for those who are not yupt and zupj.} \\
\\ 
\texttt{However, we are given that for those who are yupt and are zupj, the probability of muvq is 10\%, and for those who are not yupt and are zupj, the probability of muvq is 51\%. Since these probabilities are not the same, the path from yupt to muvq through zupj is confounded by the variable xyfo.} \\
\\ 
\texttt{However, we are given that the overall probability of xyfo is 22\%, and we can use this information to calculate the probability of muvq for those who are yupt and zupj, given that they are not xyfo. We can do this by using the law of total probability:} \\
\\ 
\texttt{P(muvq|yupt, zupj) = P(muvq|yupt, zupj, xyfo) * P(xyfo|yupt, zupj) + P(muvq|yupt, zupj, not xyfo) * P(not xyfo|yupt, zupj)} \\
\\ 
\texttt{We are given that P(muvq|yupt, zupj, xyfo) = 10\% and P(xyfo|yupt, zupj) = 44\% (since the probability of xyfo is 22\% and the probability of zupj is 44\% for those who are yupt). We are also given that P(muvq|yupt, zupj, not xyfo) = 57\% (since the probability of muvq is 57\% for those who are yupt and not zupj, and the probability of zupj is 26\% for those who are yupt and not xyfo). We can also calculate P(not xyfo|yupt, zupj) = 1 {-} P(xyfo|yupt, zupj) = 1 {-} 44\% = 56\%.} \\
\\ 
\texttt{Plugging in these values, we get:} \\
\\ 
\texttt{P(muvq|yupt, zupj) = 10\% * 44\% + 57\% * 56\% ≈ 32.28\%} \\
\\ 
\texttt{However, we are given that} \\
            }
        }
    }
    & & \\ \\

    \theutterance \stepcounter{utterance}  
    & & & \multicolumn{2}{p{0.3\linewidth}}{
        \cellcolor[rgb]{0.9,0.9,0.9}{
            \makecell[{{p{\linewidth}}}]{
                \texttt{\tiny{[GM$|$GM]}}
                \texttt{No} \\
            }
        }
    }
    & & \\ \\

    \theutterance \stepcounter{utterance}  
    & & & \multicolumn{2}{p{0.3\linewidth}}{
        \cellcolor[rgb]{0.9,0.9,0.9}{
            \makecell[{{p{\linewidth}}}]{
                \texttt{\tiny{[GM$|$GM]}}
                \texttt{yes} \\
            }
        }
    }
    & & \\ \\

    \theutterance \stepcounter{utterance}  
    & & & \multicolumn{2}{p{0.3\linewidth}}{
        \cellcolor[rgb]{0.9,0.9,0.9}{
            \makecell[{{p{\linewidth}}}]{
                \texttt{\tiny{[GM$|$GM]}}
                \texttt{game\_result = LOSE} \\
            }
        }
    }
    & & \\ \\

\end{supertabular}
}

\end{document}
